\documentclass{article}
\usepackage{graphicx} % Required for inserting images
\usepackage{amsmath,amssymb,amsthm}
\usepackage{mathdots}
\usepackage{bbm}
\usepackage[utf8x]{inputenc}
\usepackage{latexsym}
\usepackage{graphicx}
\graphicspath{{figures/}}
\usepackage{makeidx}
\usepackage{psfrag}
\usepackage{subcaption}
\usepackage{diagbox}
\usepackage{centernot}
\bibliographystyle{plain}
\newcommand{\dsp}{\displaystyle}
\newcommand{\BBM}{\begin{bmatrix}}
\newcommand{\EBM}{\end{bmatrix}}
\newcommand{\BEQ}{\begin{equation}}
\newcommand{\EEQ}{\end{equation}}
\newcommand{\BIT}{\begin{itemize}}
\newcommand{\EIT}{\end{itemize}}
\newcommand{\cf}{{\it cf.}}
\newcommand{\eg}{{\it e.g.}}
\newcommand{\ie}{{\it i.e.}}
\newcommand{\st}{\mbox{subject to}}
\newcommand{\ones}{\mathbbm{1}}
\newcommand{\reals}{\mathbb{R}}
\newcommand{\ereals}{\bar{\reals}}
\newcommand{\pereals}{(-\infty,+\infty]}
\newcommand{\naturals}{\mathbb{N}}
\newcommand{\symm}{\mathbb{S}}
\newcommand{\quadcone}{\mathbb{Q}}
\newcommand{\expcone}{K_{\exp}}
\newcommand{\complex}{\mathbb{C}}
\newcommand{\integers}{\mathbb{Z}}
\newcommand{\psimplex}{\Delta}
\newcommand{\eqbydef}{\mathrel{\stackrel{\Delta}{=}}}
\newcommand{\im}{\mathrm{i}}
\newcommand{\minimize}{\mbox{minimize}}
\newcommand{\maximize}{\mbox{maximize}}
\newcommand{\Range}{\mathcal{R}}
\newcommand{\nullspace}{\mathcal{N}}
\newcommand{\feasset}{\mathcal{F}}
\newcommand{\probdom}{\mathcal{D}}

\newcommand{\argit}[2]{#1_{#2}}
\newcommand{\xit}[1]{\argit{x}{#1}}
\newcommand{\majf}{\bar f}
\newcommand{\maj}[2]{\majf(#1\,\vert\,#2)}
\newcommand{\indicator}[1]{I_{#1}}


\DeclareMathOperator{\Span}{span}
\DeclareMathOperator{\nullity}{nullity}
\DeclareMathOperator{\rank}{rank}
\DeclareMathOperator{\Tr}{tr}
\DeclareMathOperator{\diag}{diag}
\DeclareMathOperator{\blkdiag}{blkdiag}
\DeclareMathOperator{\vect}{vec}
\DeclareMathOperator{\svect}{svec}
\DeclareMathOperator{\dom}{dom}
\DeclareMathOperator{\conv}{conv}
\DeclareMathOperator{\relint}{relint}
\DeclareMathOperator{\interior}{int}
\DeclareMathOperator{\aff}{aff}
\DeclareMathOperator{\cone}{cone}
\DeclareMathOperator{\epi}{epi}
\DeclareMathOperator{\cl}{cl}
\DeclareMathOperator*{\argmin}{argmin}
\DeclareMathOperator*{\argmax}{argmax}
\DeclareMathOperator{\sgn}{sgn}
\DeclareMathOperator{\prox}{prox}
\DeclareMathOperator{\nnz}{nnz}

\DeclareMathOperator{\sqr}{sqr}
\DeclareMathOperator{\adj}{adj}
\DeclareMathOperator{\bdiag}{bdiag}

\DeclareMathOperator{\Expect}{\mathbb{E}}
\DeclareMathOperator{\Var}{Var}

\newcommand{\E}[2][]{\Expect_{#1}\!\left[\,#2\,\right]}   % Expectation with argument
\newcommand{\CE}[2]{\E{#1 \,\middle|\, #2}}                              % Conditional expectation

\DeclareMathOperator{\Probability}{\mathbb{P}}
\newcommand{\Prob}[2][]{\Probability_{#1}\!\left[\,#2\,\right]}     % Probability with argument
\newcommand{\CProb}[2]{\Prob{#1 \,\middle|\, #2}}                       % Conditional probability

\newcommand{\tnabla}{\widetilde{\nabla}}


\newcommand{\mfile}[1]{\txtmint{#1}}
\newcommand{\mfilelabel}[2]{\txtmint{#1}}

\title{Conic Formulations for solving the Bellman Equation}
\author{Anders Hansson}
\date{August 2024}

\begin{document}

\maketitle

\section{The Bellman Equation}
%
Let $F:\mathcal D^n\times\mathcal E^m\rightarrow\mathcal D^n$ and $f:\mathcal D^n\times\mathcal E^m\rightarrow\reals$
be given functions and consider an infinite horizon optimal control problem defined as
%
\BEQ
\begin{aligned}
\minimize \;&\sum_{k=0}^{\infty}\gamma^k
f(x_k,u_k)\\
\st \;&x_{k+1}=F(x_k,u_k),\quad k\in\integers_+\\
&u_k\in U(x_k)\subseteq \mathcal D^m,\quad k\in\integers_+
\end{aligned}\label{eqn:infinite-control}
\EEQ
%
with variables $(u_0,x_1,u_1,x_2\ldots)$,
where $x_0$ is given, and where $0<\gamma\leq 1$ is a discount factor. 
Here $U$ is a set-valued function from $\mathcal D^n$. 
We define the optimal value function
$J^\star:{\mathcal D}^n\rightarrow\reals$ as the optimal value of the above
optimization problem for different values of the initial state $x_0$. 
Before we continue we will make some assumptions. It is assumed that 
$F(0,0)=0$ and that $0\in U(0)$
so that zero is a stationary point of the dynamical equation.

If there exists a positive solution $V:{\mathcal D}^n\rightarrow\reals_+$
to the {\it Bellman equation}\index{Bellman!equation}
%
\BEQ
V(x)=\min_{u\in U(x)}\left\lbrace f(x,u)+\gamma V(F(x,u))\right\rbrace,
\label{eqn:bellman}
\EEQ
%
then $V(x)=J^\star(x)$\footnote{In case the solution to the Bellman equation is not
unique only one of them is equal to the optimal cost.} 
and $u_k^\star=\mu(x_k)$ is the optimal feedback control, 
where $\mu:\mathcal D^n\rightarrow\mathcal E^m$ 
is defined as $\mu(x)=\argmin_{u \in U(x)}\left\lbrace
f(x,u)+\gamma V(F(x,u))\right\rbrace$.

We define the iterates
%
\BEQ
V_{k+1}(x)=\min_{u\in U(x)}\left\lbrace f_k(x,u)+\gamma 
V_k\left(F_k(x,u)\right)\right\rbrace
\EEQ
%
with initial value $V_0(x)=0$. The iteration process is called 
{\it value iteration} (VI).\index{value!iteration}\index{VI|see{value iteration}} 
The proof of convergence of VI to a solution of the Bellman
equation is based on
the contraction property of the {\it Bellman operator} \index{Bellman!operator}
$T:\reals^{{\mathcal D}^n}:\rightarrow\reals^{{\mathcal D}^n}$\footnote{The set
$\reals^{{\mathcal D}^n}$ is the set of all functions from $\mathcal D^n$ to 
$\reals$.}  defined as
%
\BEQ
T(V)(x)=\min_{u\in U(x)}\left\lbrace f(x,u)+\gamma V(F(x,u))\right\rbrace.\label{eqn:bellman-operator}
\EEQ
%
Contraction assumes in most cases that $\gamma<1$. 
The operator is also monotone. Our intention is now to show that the solution to the Bellman equation can
also be obtained as the solution of an LP. To this end 
we now start the VI
with a $V_0$ such that $V_1=T(V_0)\geq V_0$ for all $x\in{\mathcal D}^n$. 
One possible choice is $V_0=0$ for all $x\in\mathcal D^n$. 
We obtain 
$T(V_1)\geq T(V_0)$ from the monotonicity property and hence
$V_2=T(V_1)\geq V_1\geq V_0$. If we repeat this we obtain 
$V_k=T(V_{k-1})\geq V_{k-1} \geq V_0$. Since $V_k$ converges to the 
solution $V$ of the Bellman equation we have shown that $V$ is the 
maximum element, see \cite[p. 45]{boy+van04} of the set of functions 
$V$ that satisfy the linear 
inequalities
%
$$V(x)\leq f(x,u)+\gamma V(F(x,u)),\quad \forall (x,u)\in{\mathcal D}^n\times{\mathcal E}^m\; {\rm such\; that}\; u\in U(x).$$
%
The maximum element can be obtained by solving the LP
%
\begin{align}
\maximize\;&\sum_{x\in \mathcal D^n}c(x)V(x)\label{eqn:Bellman-LP-obj}\\
\st \;&V(x)\leq f(x,u)+\gamma V(F(x,u)),\;\forall (x,u)\in {\mathcal D}^n\times{\mathcal E}^m
\; {\rm such\; that}\; u\in U(x),
\label{eqn:Bellman-LP-con}
\end{align}
%
where $c(x)>0$ is arbitrary. 
This follows from the dual characterization of maximum elements, see
\cite[p. 54]{boy+van04}. The optimization
variable is $V(x)$ for all values of $x$. In case $\mathcal D^n$ is not
finite as is the case when $\mathcal D^n=\reals^n$, we can replace
the summation with an integration, see Appendix.  
%
\section{LP Formulation}
%
Now consider the special case when $f(x,u)=s^Tx+r^Tu$ and 
$F(x,u)=Ax+Bu$ where $x\in\reals^n$, $u\in\reals^m$ and the vectors
$r$ and $s$ and the matrices $A$ and $B$ are real-valued and of 
compatible dimensions. Also let $u\in U(x)$ be defined by $|u|\leq Ex$ for a
real-valued matrix of compatible dimensions. We then obtain the LP
%
\begin{align*}
\maximize\;&\int_{x\in \reals^n}c(x)V(x)dx\\
\st \;&V(x)\leq s^Tx+r^Tu+\gamma V(Ax+Bu),\;\forall (x,u)\in\reals^{n+m}\;
{\rm such\; that}\; |u|\leq Ex.
\end{align*}
%
If we guess that $V(x)=p^Tx$ for some vector $p\in\reals^n$, then we obtain
the constraint
%
$$p^Tx\leq s^Tx+r^Tu+\gamma p^T(Ax+Bu),\;\forall (x,u)\in\reals^{n+m}\;
{\rm such\; that}\; |u|\leq Ex.$$
%
which is equivalent to 
%
$$p\leq s+\gamma A^Tp-E^T|r+\gamma B^Tp|$$
%
We also have that $\int_{x\in \reals^n}c(x)p^Txdx=p^T\mu$ if we take 
$c(x)$ to be any
probability density function that is positive and has first moments all
equal to $\mu\in\reals_{++}^n$.\footnote{See Appendix for a motivation for this choice of first moment.}
We now conclude that the originally infinite dimensional LP is 
equivalent to the finite dimensional LP
%
\begin{align*}
\maximize\;&\mu^Tp\\
\st \;&p\leq s+\gamma A^Tp-E^T\alpha\\
&-\alpha\leq r+\gamma B^Tp\leq \alpha
\end{align*}
%
with variable $p$. Notice that the solution $p$ is the maximum element that satisfies
the constraints of the LP. This follows from the dual characterization of maximum 
elements. 
%
\section{Conic Optimal Control Problem}
%
We will in this section look at a generalization of the previous control problem. 
We let the incremental cost be the same and the dynamics, but the set $U(x)$ is defined
as
%
$U(x)=\{u|(x,u)\in \cal P\}$
%
where $\cal P$ is a proper convex cone. Let $\mathcal X=\{x|\exists u: (x,u)\in\mathcal P\}$.
We will also assume that there exists a function
$\phi:{\mathcal E}^m\rightarrow{\mathcal D}^n$ such that
%
$$\min_{u\in U(x)}\lambda^Tu=\phi(\lambda)^Tx$$
%
This implies that we can use  $V(x)=p^Tx$ also for our more general 
problem. Hence we will now have the infinite-dimensional LP formulation
%
\begin{align*}
\maximize\;&\int_{x\in \mathcal X}c(x)p^Txdx\\
\st \;&p^Tx\leq s^Tx+r^Tu+\gamma p^T(Ax+Bu),\;\forall (x,u)\in \mathcal P
\end{align*}
%
With $c$ as any positive probability density function defined on $\mathcal X$ and with 
first moment $\mu\in\reals_{++}^n$ we obtain objective function $\mu^Tp$. We rewrite the
constraint as
%
$$(s-p+\gamma A^Tp)^Tx+(r+\gamma B^Tp)u\geq 0,\;\forall (x,u)\in \mathcal P$$
%
and by the definition of the dual cone $\mathcal P^\star$ we have that the 
constraint is equivalent to $(s-p+A^Tp,r+B^Tp)\in\mathcal P^\star$, and therefore
we obtain the following conic optimization problem for determining $p$
%
\begin{align*}
\maximize\;&\mu^Tp\\
\st \;&(s-p+\gamma A^Tp,r+\gamma B^Tp)\in\mathcal P^\star
\end{align*}
%
with variable $p$.
%
\section{SDP Formulation}
%
We now consider the LQ control problem which is obtained by changing the 
incremental cost to 
%
$$f(x,u)=\BBM x\\u\EBM^TQ\BBM x\\u\EBM$$
%
for a matrix $Q\in\symm_+^{m+n}$. We still let 
%
$$F(x,u)=\BBM A&B\EBM\BBM x\\u\EBM$$ 
%
The guess is that $V(x)=x^TPx$ for some $P\in\symm^n$. The infinite dimensional
LP will then be
%
\begin{align*}
\maximize\;&\int_{x\in \reals^n}c(x)x^TPxdx\\
\st \;&x^TPx\leq \BBM x\\u\EBM^TQ\BBM x\\u\EBM+\gamma 
\BBM x\\u\EBM^T\BBM A&B\EBM^TP\BBM A&B\EBM\BBM x\\u\EBM,\;\forall (x,u)\in\reals^{n+m}.
\end{align*}
%
with variable $P$. Also here $c$ is a positive probability density function, but
this time we will let it have zero mean and covariance $\Sigma\in\symm_{++}^n$.\footnote{See Appendix for a motivation for this choice of moments.}
We then realize that the objective function becomes $\Tr(\Sigma P)$ and that the
constraints are equivalent to the LMI
%
$$\BBM P&0\\0&0\EBM\preceq Q+\gamma \BBM A&B\EBM^TP\BBM A&B\EBM$$
%
Hence we obtain the SDP
%
\begin{align*}
\maximize\;&\Tr(\Sigma P)\\
\st \;&\BBM P&0\\0&0\EBM\preceq Q+\gamma \BBM A&B\EBM^TP\BBM A&B\EBM
\end{align*}
%
with variable $P$. Notice that the solution $P$ is the maximum element that satisfies
the constraints of the SDP. 
%
\section{Conclusions}
%
It is interesting to note that the property of the value function 
$V$ solving the Bellman equation being the maximum element of a set 
described by an infinite number of affine constraints is translated into the property
that the vector or matrix that parameterizes the values function is the
maximum element of a set that is a finite-dimensional translated cone. 
%
\section{Appendix: Dual Characterization}
%
Let $\mathcal X\subseteq \reals^n$. 
Suppose $V:\mathcal X\rightarrow\reals_+$ is a maximum element of a set $S\subseteq \reals_+^{\mathcal X}$, i.e.
$V(x)\geq U(x)$ for all $x\in\mathcal X$ and all $U\in S$.
Consider any $c:\mathcal X\rightarrow \reals_{++}$ 
such that $0<\langle c,V\rangle=\int_{\mathcal X}c(x)V(x)dx<\infty$. 
Then let $U\in S$, $U\neq V$. Since $V$ is a maximum
element of $S$ we have that $V(x)-U(x)\geq 0$ for all $x\in\mathcal X$. From this it follows that 
$\langle c,V-U\rangle>0$ or equivalently that $\langle c,V\rangle > \langle c,U\rangle$. Since
$U$ is an arbitrary element in $S$ not equal to $V$, this shows that $V$ is the unique maximizer of 
$\langle c,U\rangle$ over $U\in S$, modulo a set of points of measure zero. 
We may take $c$ to be probability density functions. Then the condition 
$0<\langle c,V\rangle<\infty$ is equivalent to  $0<\E{V(x)}<\infty$.  If we know that 
$V$ is linear in $x$ we see that this condition is satisfied if the first moment is finite and
nonzero, and if
$V$ is quadratic in $x$ we see that this condition is satisfied if also the second moment is finite and
nonzero. However, we need to impose more conditions. We are most often interested in solutions
$V$ of the Bellman equation which are such that $V(x)>0$ whenever $x\neq 0$. 
If $V$ is linear in $x$, i.e. there is a vector $p\in\reals^n$ such that $V(x)=p^Tx$, then it must be
that $p\in\reals_{++}^n$. We also have that $\E{V(x)}=p^T\mu$, where $\E{x}=\mu$. In case we do not 
ensure that $\mu\in\reals_{++}^n$, it will be that $p_i=0$ for some $i$ could be optimal. In order for 
this to never happen we will impose the constraint that $\mu\in\reals_{++}^n$. Similarly if 
$V$ is quadratic in $x$ we will also impose the constraint that  $\E{xx^T}=\Sigma\in
\symm_{++}^n$. 
%
\bibliography{references}
\end{document}
